\documentclass[11pt]{article}
\usepackage[utf8]{inputenc}
\usepackage[spanish]{babel}
\usepackage{amsmath,amssymb,amsthm}
\usepackage{graphicx}
\usepackage[margin=2.5cm]{geometry}
\usepackage{enumitem}

\title{Semana 1 -- Ejercicios}
\author{Tutorías -- Introducción al Cálculo}
\date{03/08--07/08/2026}

\begin{document}
\maketitle

\section{Ejercicios de la semana}

\begin{enumerate}[label=\arabic*.]

\item Determine el grado de los siguientes polinomios, si es que tienen:
\begin{enumerate}[label=(\alph*)]
\item $P(x) = 5x^3 - 2x + 7$
\item $P(x) = 0 \cdot x^4 + 0 \cdot x + 0$
\item $P(x) = -8$
\end{enumerate}
(Respuestas: 3; sin grado; 0.)

\item Sean $P(x) = 2x^3 - x^2 + 4$ y $Q(x) = x^2 - 3x + 1$. Calcule $P(x)+Q(x)$,
$P(x)-Q(x)$ y $P(x)\cdot Q(x)$.

\item Verifique, usando la definición de raíz, si $\alpha = 3$ es raíz de
$P(x) = x^3 - 4x^2 + x + 6$. (Respuesta: no lo es.)

\item Use el Teorema del Resto para calcular el resto de dividir
$P(x) = 2x^4 - 3x^3 + x - 1$ por $(x-1)$, sin desarrollar la división
completa. (Respuesta: $R=-1$.)

\item Use el esquema de coeficientes indeterminados para dividir
$P(x) = x^3 + 2x^2 - 5x - 6$ por $(x+3)$. Indique el cuociente y el resto, y
concluya si $-3$ es o no raíz de $P$. (Respuesta: $Q(x)=x^2-x-2$, $R=0$;
$-3$ sí es raíz.)

\item Divida $P(x) = 4x^3 - x^2 + 2x - 5$ por $(x-\tfrac12)$ usando el
esquema de coeficientes indeterminados.

\item Aplique el algoritmo general de división (Teorema 7.6) para dividir
$P(x) = x^4 + x^3 - 3x^2 + x - 2$ por $S(x) = x^2 + 2x - 1$. Encuentre $Q$ y
$R$, y verifique que $Q(x)S(x)+R(x) = P(x)$.

\item Divida $P(x) = 2x^5 - 3x^2 + 1$ por $S(x) = x^3 - x + 1$.

\item ¿Es posible que al dividir un polinomio $P$ de grado $4$ por un
polinomio $S$ de grado $5$, el cuociente $Q$ no sea el polinomio $0$?
Justifique usando el Teorema 7.6. (Respuesta: no.)

\item Sin hacer la división completa, determine si $(x+2)$ divide a
$P(x) = x^5 + 3x^4 - x^2 + 4x + 12$. (Respuesta: sí, ya que $P(-2)=0$.)

\item \textbf{(Nivel interrogación)} Sea $P(x)$ un polinomio de grado $n\ge1$ y $S(x)$ un
polinomio de grado $1$. Usando el Teorema 7.6 (algoritmo general de división), demuestre
rigurosamente que el resto $R(x)$ de dividir $P(x)$ por $S(x)$ es necesariamente un polinomio
\emph{constante} (es decir, de grado $0$).

\item \textbf{(Nivel interrogación)} Sea $P(x) = x^4 - 2x^3 + kx^2 - 3x + 5$, con $k\in\mathbb{R}$.
Sabiendo que el resto de dividir $P(x)$ por $(x+1)$ es $12$, determine el valor de $k$ (use el
Teorema del Resto, sin hacer la división completa).

\item \textbf{(Nivel interrogación)} Sea $P(x)$ un polinomio de grado $n$ y $S(x)$ un polinomio de
grado $m$ con $m>n$. Demuestre, a partir del Teorema 7.6, que el cuociente de dividir $P(x)$ por
$S(x)$ es el polinomio $0$ y que el resto es el propio $P(x)$.

\end{enumerate}

\end{document}
